\documentclass[11pt,a4paper]{article}

\usepackage[utf8]{inputenc}
\usepackage[T1]{fontenc}
\usepackage{lmodern}
\usepackage[margin=2.5cm]{geometry}
\usepackage{booktabs}
\usepackage{tabularx}
\usepackage{xcolor}
\usepackage{hyperref}
\usepackage{listings}
\usepackage{enumitem}
\usepackage{float}
\usepackage{caption}

\definecolor{codeblue}{HTML}{2060A0}
\definecolor{codegray}{HTML}{555555}
\definecolor{codegreen}{HTML}{208030}
\definecolor{backgray}{HTML}{F5F5F5}
\hypersetup{colorlinks=true, linkcolor=codeblue, citecolor=codeblue, urlcolor=codeblue}

\lstdefinestyle{plain}{
  backgroundcolor=\color{backgray},
  basicstyle=\small\ttfamily,
  breaklines=true,
  frame=single,
  rulecolor=\color{codegray!30},
  showstringspaces=false,
}
\lstset{style=plain}

\title{\textbf{Delphi}\\[4pt] \large A Multi-Agent World-Simulation Engine \\[2pt] \normalsize Extended for Financial-Market Scenario Prediction --- Documentation}
\author{Vito Ciciretti}
\date{July 2026}

\begin{document}
\maketitle
\tableofcontents

\section{Overview}

Delphi builds a high-fidelity parallel digital world from real-world seed
material --- a news article, a policy draft, a financial signal, or a story to
re-imagine. Inside that world, agents with independent personalities, long-term
memory, and behavioural logic interact and evolve. Variables can be injected
from a god's-eye view to watch how the resulting dynamics play out --- rehearsing
outcomes in a sandbox before they happen for real.

Given seed material and a natural-language prediction question, Delphi returns a
detailed report and an interactive digital world that remains explorable after
the simulation ends.

\subsection{Intended use}

\begin{itemize}[nosep]
  \item \textbf{Macro} --- a zero-risk rehearsal lab for decision-makers: test
    policies, messaging, and market moves before committing to them.
  \item \textbf{Micro} --- a creative sandbox: deduce alternate story endings or
    explore ``what if'' scenarios.
\end{itemize}

\section{Provenance}
\label{sec:provenance}

Delphi is a \textbf{derivative work}. Its simulation core is
\href{https://github.com/666ghj/MiroFish}{MiroFish}, an AGPL-3.0 multi-agent
prediction engine, and its underlying agent-interaction substrate is
\href{https://github.com/camel-ai/oasis}{OASIS} (Open Agent Social Interaction
Simulations) from the CAMEL-AI team. Delphi preserves the AGPL-3.0 license end
to end, as required. The full list of modifications is tracked in
\texttt{NOTICE.md} in the repository root.

\textbf{What is original to Delphi} is described in Section~\ref{sec:presets}:
a config-driven scenario/domain preset layer that generalises the engine from a
single fixed scope (social-media public opinion) to arbitrary domains, including
financial markets.

\section{Workflow}

\begin{enumerate}[nosep]
  \item \textbf{Graph building} --- seed extraction, individual/collective
    memory injection, GraphRAG construction.
  \item \textbf{Environment setup} --- entity relationship extraction, persona
    generation, agent configuration injection.
  \item \textbf{Simulation} --- parallel multi-channel simulation, auto-parsed
    prediction requirements, dynamic temporal memory updates.
  \item \textbf{Report generation} --- a report agent with a tool set for deep
    interaction with the post-simulation world.
  \item \textbf{Deep interaction} --- chat with any agent in the simulated
    world, or with the report agent itself.
\end{enumerate}

\section{Scenario / Domain Presets}
\label{sec:presets}

The upstream engine hard-wired its assumptions (rhythm, channels, stance
vocabulary, prompt framing) to a single scope: social-media public-opinion
simulation, on a fixed China-timezone activity rhythm. Delphi externalises all of
that into \textbf{scenario presets} --- config-driven JSON descriptions of what
kind of world a simulation runs in --- so the same engine can model social media,
financial markets, organisational decisions, or narrative worlds without
touching engine code.

\subsection{What a preset controls}

\begin{table}[H]
\centering
\begin{tabularx}{\textwidth}{@{}l X@{}}
\toprule
\textbf{Field} & \textbf{Meaning} \\
\midrule
\texttt{id} & Unique preset identifier (used in the API / \texttt{SCENARIO\_DEFAULT}). \\
\texttt{name} / \texttt{description} & Human-facing label and blurb shown in the UI. \\
\texttt{domain} & High-level grouping: \texttt{social\_media}, \texttt{market}, \texttt{organization}, \texttt{narrative}, \texttt{custom}. \\
\texttt{activity\_rhythm} & When agents are active over 24h (peak / work / morning / off-peak hours and multipliers) plus a timezone label. Replaces the old fixed China-timezone rhythm. \\
\texttt{default\_total\_hours} / \texttt{default\_minutes\_per\_round} & Default simulation horizon and time granularity. \\
\texttt{channels} & Communication channels agents use; each maps onto an \texttt{engine\_platform} (\texttt{twitter} or \texttt{reddit} for OASIS) with recommender weights, a viral threshold, and echo-chamber strength. \\
\texttt{stances} / \texttt{default\_stance} & Domain-appropriate stance vocabulary (e.g.\ bullish/bearish for markets, advocate/skeptic for organisations). \\
\texttt{prompt\_framing} & Natural-language framing injected into LLM prompts so generated events, personas, and behaviours fit the domain. \\
\texttt{tags} & Free-form labels for filtering and display. \\
\bottomrule
\end{tabularx}
\caption{Fields exposed by a scenario preset.}
\end{table}

\textbf{Engine note.} The underlying simulation still runs on OASIS social
platforms, so every channel maps onto \texttt{twitter} or \texttt{reddit}. A
preset changes the framing, rhythm, and dynamics of a run; the \texttt{engine}
field on each channel leaves room for alternative simulation backends later.

\subsection{Built-in presets}

\begin{table}[H]
\centering
\begin{tabularx}{\textwidth}{@{}l l X@{}}
\toprule
\textbf{id} & \textbf{Domain} & \textbf{Highlights} \\
\midrule
\texttt{social\_media} & social\_media & \textbf{Default.} Reproduces the original engine behaviour: China-timezone rhythm, 72h, Twitter + Reddit. \\
\texttt{global\_social\_media} & social\_media & Flattened 24/7 rhythm for worldwide, cross-timezone events. \\
\texttt{financial\_market} & market & Market-hours rhythm, 48h in 30-minute rounds, fast news feed + investor forum, bullish/bearish stances. \\
\texttt{organization} & organization & Business-hours rhythm, 5-day horizon, announcements + team discussion, advocate/skeptic/blocker stances. \\
\texttt{creative\_narrative} & narrative & No diurnal cycle, week-long story-time horizon, in-character framing for alternate endings. \\
\bottomrule
\end{tabularx}
\caption{Presets shipped with Delphi.}
\end{table}

\subsection{Selecting and adding scenarios}

Pass \texttt{scenario\_id} when creating a simulation:

\begin{lstlisting}
POST /api/simulation/create
{
  "project_id": "proj_xxxx",
  "scenario_id": "financial_market"
}
\end{lstlisting}

If omitted, \texttt{Config.SCENARIO\_DEFAULT} (env \texttt{SCENARIO\_DEFAULT},
default \texttt{social\_media}) is used. Which platforms run is inferred from the
preset's channels unless \texttt{enable\_twitter} / \texttt{enable\_reddit} is
passed explicitly. List everything available with
\texttt{GET /api/simulation/scenarios}.

Adding a new domain requires no code: drop a \texttt{<id>.json} file into
\texttt{SCENARIO\_PRESETS\_DIR} (default \texttt{backend/uploads/scenarios}),
following the schema of the built-ins in
\texttt{backend/app/scenarios/presets/}. A user preset with the same \texttt{id}
as a built-in overrides it, so a shipped scenario can be retuned without editing
source. Presets are validated on load; malformed files raise a clear
\texttt{ScenarioLoadError} at startup rather than failing silently at runtime.

Minimal example:

\begin{lstlisting}
{
  "id": "product_launch",
  "name": "Product Launch Buzz",
  "domain": "social_media",
  "activity_rhythm": { "timezone": "UTC",
                        "peak_hours": [12, 13, 18, 19, 20] },
  "default_total_hours": 96,
  "channels": [
    { "id": "social", "engine_platform": "twitter",
      "viral_threshold": 8 },
    { "id": "forum",  "engine_platform": "reddit" }
  ],
  "stances": ["excited", "critical", "neutral", "observer"],
  "default_stance": "neutral",
  "prompt_framing": "Simulate public reaction to a new product launch..."
}
\end{lstlisting}

\section{Setup}

\subsection{Prerequisites}

\begin{table}[H]
\centering
\begin{tabularx}{\textwidth}{@{}l l X@{}}
\toprule
\textbf{Tool} & \textbf{Version} & \textbf{Description} \\
\midrule
Node.js & 18+ & Frontend runtime, includes npm \\
Python & $\geq$3.11, $\leq$3.12 & Backend runtime \\
uv & latest & Python package manager \\
\bottomrule
\end{tabularx}
\end{table}

\subsection{Source deployment}

\begin{lstlisting}
cp .env.example .env
# fill in LLM_API_KEY / LLM_BASE_URL / LLM_MODEL_NAME
# and ZEP_API_KEY (memory graph; free tier from app.getzep.com)

npm run setup:all      # install root + frontend + backend deps
npm run dev            # start frontend (:3000) + backend (:5001)
\end{lstlisting}

\subsection{Docker deployment}

\begin{lstlisting}
cp .env.example .env
docker compose up -d --build
\end{lstlisting}

\section{License \& Attribution}

Delphi is licensed under the GNU Affero General Public License v3.0 --- see
\texttt{LICENSE}. As described in Section~\ref{sec:provenance}, it is a
derivative of MiroFish (AGPL-3.0) and depends on OASIS (CAMEL-AI). All original
copyright and license notices are retained; trademarks, logos, and project names
referenced from the upstream project remain the property of their respective
owners. The complete, itemised list of modifications is in \texttt{NOTICE.md}.

\end{document}
